\documentclass[11pt,letterpaper]{article}
\usepackage[T1]{fontenc}
\usepackage[utf8]{inputenc}
\usepackage{amsmath,amsthm}
\usepackage{newpxtext,newpxmath}
\usepackage{microtype,mathtools}
\usepackage[margin=1in,top=0.85in,bottom=0.9in]{geometry}
\usepackage{booktabs,longtable,array,enumitem}
\usepackage[dvipsnames]{xcolor}
\usepackage[most]{tcolorbox}
\usepackage{fancyhdr}
\usepackage{hyperref}
\definecolor{InkGreen}{HTML}{3F513C}
\definecolor{PaleGreen}{HTML}{F2F5EF}
\definecolor{RuleGreen}{HTML}{C3CEBC}
\hypersetup{colorlinks=true,linkcolor=InkGreen,citecolor=InkGreen,urlcolor=InkGreen,
 pdftitle={Phase-shift classification for integer tetration},
 pdfauthor={ChatGPT; prepared for Vladimir Reshetnikov}}
\pagestyle{fancy}
\fancyhf{}
\fancyhead[L]{\small\color{InkGreen}Integer tetration: phase-shift classification}
\fancyhead[R]{\small Research draft}
\fancyfoot[C]{\thepage}
\setlength{\headheight}{14pt}
\setlength{\parskip}{0.3em}
\setlength{\emergencystretch}{2em}
\numberwithin{equation}{section}
\newtheorem{theorem}{Theorem}[section]
\newtheorem{lemma}[theorem]{Lemma}
\newtheorem{proposition}[theorem]{Proposition}
\newtheorem{corollary}[theorem]{Corollary}
\theoremstyle{definition}
\newtheorem{definition}[theorem]{Definition}
\newtheorem{remark}[theorem]{Remark}
\newcommand{\N}{\mathbb N}
\newcommand{\F}{\mathbb F}
\newcommand{\Z}{\mathbb Z}
\newcommand{\APS}{\operatorname{APS}}
\newcommand{\MPS}{\operatorname{MPS}}
\newcommand{\ord}{\operatorname{ord}}
\newcommand{\vp}[2]{v_{#1}(#2)}
\newcommand{\dd}{\mathrm d}
\newcommand{\Eset}{\mathcal E}
\newcommand{\Uset}{\mathcal U}
\newtcolorbox{resultbox}{colback=PaleGreen,colframe=RuleGreen,boxrule=.6pt,
 arc=1mm,left=3mm,right=3mm,top=2mm,bottom=2mm,breakable}
\title{\color{InkGreen}\LARGE Phase-shift classification\par for integer tetration\par
\vspace{0.5em}\large A proof of the 21-value conjecture, its even-base refinement,\par
and the related 28-value modular classification}
\author{ChatGPT\\\small Prepared for Vladimir Reshetnikov}
\date{19 September 2026}
\begin{document}
\maketitle
\thispagestyle{empty}
\begin{abstract}
Integer power towers become extraordinarily large, but the first decimal digit
that changes between successive heights has a small and rigid eventual dynamics.
Marco Rip\`a's phase-shift conjecture, stated in the appendix of his 2025 paper
and in OEIS A376842, predicts exactly 21 possible phase words when the cycle is
anchored at the first height of permanently constant congruence speed.
We give a self-contained elementary proof under the height-one convention
explicitly recorded in OEIS A371074. The proof separates the 2-adic and 5-adic
valuations of consecutive tower differences. Their competition determines whether
the phase is always 5, follows multiplication in $\F_5^\times$, or belongs to a
balanced family $a=1+10^v c$. The balanced family accounts for precisely four
additional words. We also prove the conjectured 14-word restriction for even
bases, the 28-word classification at the alternative anchor of OEIS A376446,
and natural density $1/900$ for each of the four words arising from balanced
valuations. Explicit onset formulas, an unbounded sharp-onset family, a
polynomial-time computation argument, and independently cross-checked exact
integer programs accompany the proofs. A source audit identifies a height-zero
indexing ambiguity and an apparent erroneous value at base 21; neither is used
as a substitute for proving the conjecture.
\end{abstract}
\begin{resultbox}
\textbf{Status of this document.} This is a proof-bearing research draft, not a
refereed publication or a formally verified development. The universal claims
below rest on the written arguments, not on a finite search. The consulted
sources still state the target as a conjecture; no earlier proof of the complete
classification was located in the targeted literature check. That search does
not establish priority. Existing congruence-speed results are explicitly
credited and rederived where needed.
\end{resultbox}
\clearpage
\begingroup
\setlength{\parskip}{0pt}
\small
\tableofcontents
\endgroup
\clearpage

\section{The selected problem and its provenance}

Fix an integer base $a>1$ and define the right-associated power tower
\begin{equation}\label{eq:tower}
 T_0=1,\qquad T_{b+1}=a^{T_b}\quad(b\geq0).
\end{equation}
Thus $T_b=a\uparrow\uparrow b$ for positive integer $b$. This article concerns
integer tetration and modular arithmetic, not an analytic interpolation in the
height. Higher up-arrow expressions often supply enormous tower heights, but
nothing in the proof requires writing those heights or towers in decimal.

The selected problem is an explicit, finite-valued conjecture rather than a
newly invented question. Rip\`a asked a 23-value version on MathOverflow in
October 2024 \cite{mo}. The appendix of the 2025 paper gives the sharper
21-value version and an additional 14-value restriction for even bases
\cite{ripa2025}. The OEIS A376842 entry, revision 39 dated 30 March 2026,
still labels the 21-value statement a conjecture \cite{oeis842}. We address
that formulation. The older 23-value set contains two extra rotations,
$1397$ and $1793$, which do not occur at the intended anchor.

The general stabilization of trailing digits and formulas for their rate of
stabilization are established topics: see Rip\`a's 2020 and 2021 papers and the
2022 paper with Onnis \cite{ripa2020,ripa2021,ripaonnis}. Those results should
not be confused with the phase-word classification. The present proof derives
the particular valuation identities it uses, then supplies the missing
classification and its anchoring argument.

\subsection{The observable}
Put
\begin{equation}\label{eq:difference}
 D_b=T_{b+1}-T_b>0\quad(b\geq0),\qquad
 k_b=\min\{v_2(D_b),v_5(D_b)\}\quad(b\geq1).
\end{equation}
For a nonzero integer $n$, $v_p(n)$ is the largest $e\geq0$ with $p^e\mid n$.
We also write $v_{10}(n)=\min\{v_2(n),v_5(n)\}$, although 10 is not prime.
The number $k_b$ counts the trailing decimal positions in which $T_b$ and
$T_{b+1}$ agree, allowing leading zeroes. Their first differing position is
$k_b+1$, counted from the right.

\begin{definition}[Phase and speed]\label{def:phase}
For $b\geq1$, define the phase digit by
\begin{equation}\label{eq:phase}
 s_b\equiv-\frac{D_b}{10^{k_b}}\pmod {10},\qquad 1\leq s_b\leq9.
\end{equation}
Define the congruence speeds by
\begin{equation}\label{eq:speed}
 V_1=k_1,\qquad V_b=k_b-k_{b-1}\quad(b\geq2).
\end{equation}
When these speeds are eventually constant, write $v$ for the eventual value and
\begin{equation}\label{eq:onsetdef}
 B=\min\{b\geq1:V_j=v\text{ for every }j\geq b\}.
\end{equation}
The word $\APS(a)$ is the four digits $s_Bs_{B+1}s_{B+2}s_{B+3}$, reduced to two
digits when the two halves agree, and then to one digit when those two digits
agree. Words are written as integers, as in the OEIS convention.
\end{definition}

The negative sign in \eqref{eq:phase} is intentional: the original digit minus
the corresponding next-height digit is being measured. Since the lower $k_b$
digits agree, reducing $(T_b-T_{b+1})/10^{k_b}$ modulo 10 gives exactly that
digit difference, including any borrow automatically.

\subsection{A convention that changes the answer}\label{sec:convention}
The convention $V_1=k_1$ is not optional bookkeeping. It is explicitly stated
in OEIS A371074: the number of matching trailing digits of $a$ and $a^a$ is the
congruence speed at height one, with leading zeroes included \cite{oeis1074}.
It also makes $k_b=\sum_{j=1}^b V_j$, as in Equation (1) of
Rip\`a--Onnis \cite{ripaonnis}.

Although $T_0=1$ is useful for algebra, we \emph{do not} define
$V_1=k_1-v_{10}(a-1)$. That alternative shifts the starting point for some
cycles. In particular it would replace $\APS(901)=19$ by $91$.
Section~\ref{sec:audit} records the ambiguity in the source formulations and
checks it against their explicit examples. Below, $k_0^*=0$ will mean only the
artificial initial total used in \eqref{eq:speed}; it is not an assertion about
$v_{10}(D_0)$.

\subsection{Examples before the proof}
For $a=3$, the first towers are $3,27,3^{27},\ldots$. One obtains
$k_b=b-1$, $V_1=0$, $V_b=1$ for $b\geq2$, and $B=2$. The phases from height two
are $4,6,4,6,\ldots$, so $\APS(3)=46$.

For $a=57$, $k_b=3b$ and $B=1$. The successive phases are
$2,6,8,4,2,6,8,4,\ldots$. Thus a three-digit stabilization speed coexists with
a four-step phase period: $\APS(57)=2684$. The speed and the phase period are
different invariants.

\section{Main classification}
Define the following sets of words:
\begin{align}
 \Eset={}&\{2,4,6,8,28,46,64,82,\notag\\[-2pt]
 &\hspace{1.2em}2486,2684,4268,4862,6248,6842,8426,8624\},\label{eq:Eset}\\
 \Uset={}&\{9,19,3971,7931\}.\label{eq:Uset}
\end{align}

\begin{theorem}[The 21-value classification]\label{thm:main}
Let $a>1$ be an integer not divisible by 10, and use
Definition~\ref{def:phase}. The speeds become permanently constant at a finite
height $B$, with $v>0$. The phases are periodic from that very height, with
least period dividing 4, and
\begin{equation}\label{eq:mainset}
 \APS(a)\in \Eset\cup\{5\}\cup\Uset.
\end{equation}
Every one of these 21 words occurs. More precisely, exactly one of the
following alternatives holds at every height $b\geq B$:
\begin{enumerate}[label=\textup{(\roman*)},leftmargin=2em,itemsep=2pt]
\item $v_2(D_b)<v_5(D_b)$, and $s_b=5$.
\item $v_5(D_b)<v_2(D_b)$, and $s_b$ is the even lift of a geometric
      progression in $\F_5^\times$.
\item $a=1+10^v c$ with $v\geq2$ and $\gcd(c,10)=1$. In this case
      $B=2$, $k_b=(b+1)v$, and
      \begin{equation}\label{eq:balancedphase}
      s_b\equiv-c^{b+1}\pmod {10}\qquad(b\geq1).
      \end{equation}
\end{enumerate}
In alternative \textup{(iii)}, residues $c\equiv1,3,7,9\pmod {10}$ give,
respectively, the words $9,3971,7931,19$.
\end{theorem}

\begin{theorem}[The even-base refinement]\label{thm:even}
For even $a>1$ with $5\nmid a$, the exact range is
\begin{equation}\label{eq:evenset}
 \{2,4,6,8,28,46,64,82,2486,4268,4862,6248,6842,8426\}.
\end{equation}
In particular, the two even-digit words $2684$ and $8624$ occur for odd bases
but never for even bases at the minimal-speed anchor.
\end{theorem}

The important qualification is ``from that very height.'' Merely proving that
the phases eventually become periodic would not settle the conjecture:
$\APS(a)$ is anchored at the \emph{first} permanently constant speed. The
valuation comparison in Section~\ref{sec:onset} closes that gap.

\section{Elementary valuation tools}

The basic identity connecting adjacent differences is
\begin{equation}\label{eq:diffrec}
 D_b=T_b\bigl(a^{D_{b-1}}-1\bigr)\quad(b\geq1).
\end{equation}
Indeed, $T_{b+1}/T_b=a^{T_b-T_{b-1}}=a^{D_{b-1}}$. This identity is an equality
of ordinary integers; no infinite tower or $p$-adic analytic continuation is
being assumed.

\begin{lemma}[The 5-adic lifting calculation]\label{lem:five}
If $x=1+5^r u$ with $r\geq1$ and $5\nmid u$, then for every positive integer $n$,
\begin{equation}\label{eq:fiveLTE}
 v_5(x^n-1)=r+v_5(n),\qquad
 \frac{x^n-1}{5^{r+v_5(n)}}\equiv
 u\frac{n}{5^{v_5(n)}}\pmod5.
\end{equation}
\end{lemma}
\begin{proof}
For an exponent $w$ coprime to 5, the binomial theorem gives
$(1+5^r u)^w\equiv1+w5^r u\pmod {5^{r+1}}$. Raising a number
$1+5^j z$ to its fifth power gives $1+5^{j+1}z'$ with $z'\equiv z\pmod5$:
the terms of binomial degree at least two are divisible by $5^{j+2}$.
Write $n=5^e w$, apply this fifth-power step $e$ times, and then the coprime
exponent calculation. The displayed residue is nonzero, proving the exact
valuation as well as the congruence.
\end{proof}

\begin{lemma}[The 2-adic lifting calculation]\label{lem:two}
For odd $a>1$ put
\begin{equation}\label{eq:At}
 A=v_2(a-1),\qquad t=v_2(a-1)+v_2(a+1)-1
                  =\max\{v_2(a-1),v_2(a+1)\}.
\end{equation}
Then $t\geq2$. For positive even $n$,
\begin{equation}\label{eq:twoLTE}
 v_2(a^n-1)=t+v_2(n).
\end{equation}
\end{lemma}
\begin{proof}
Exactly one of $a-1,a+1$ is divisible by 4, and the other has valuation one.
For odd $w$, the factorizations of $a^w-1$ and $a^w+1$ show that their
2-adic valuations equal those of $a-1$ and $a+1$, respectively. If
$n=2^e w$ with $e\geq1$, factor the successive differences of squares.
After the first factorization, each extra factor is $a^{2^j w}+1$ with
$j\geq1$, which is congruent to 2 modulo 8. Hence the valuation is
$v_2(a-1)+v_2(a+1)+(e-1)=t+e$.
\end{proof}

For $5\nmid a$, set
\begin{equation}\label{eq:or}
 o=\ord_5(a)\in\{1,2,4\},\qquad r=v_5(a^o-1)\geq1.
\end{equation}
The value of $r$ is $v_5(a-1)$, $v_5(a+1)$, or $v_5(a^2+1)$ according as
$o=1,2,4$. In the latter two cases the other factor of $a^o-1$ is a 5-adic unit.
Whenever $o\mid n$, Lemma~\ref{lem:five} yields
\begin{equation}\label{eq:orderLTE}
 v_5(a^n-1)=r+v_5(n).
\end{equation}
If $o\nmid n$, this valuation is zero.

\section{Exact local valuations of tower differences}\label{sec:valuations}
Write $x_b=v_2(D_b)$ and $y_b=v_5(D_b)$.

\begin{proposition}[Odd bases coprime to 5]\label{prop:oddprofiles}
For odd $a>1$ with $5\nmid a$, and every $b\geq1$,
\begin{equation}\label{eq:oddprofiles}
 x_b=A+bt,\qquad y_b=(b+\varepsilon)r,
\end{equation}
where
\begin{equation}\label{eq:epsilonodd}
\varepsilon=\begin{cases}
 1,&o=1,\\
 0,&o=2,\\
 0,&o=4\text{ and }a\equiv1\pmod4,\\
 -1,&o=4\text{ and }a\equiv3\pmod4.
\end{cases}
\end{equation}
In particular, $k_b=\min\{A+bt,(b+\varepsilon)r\}$ exactly, not just
asymptotically.
\end{proposition}
\begin{proof}
All $D_b$ are even and $T_b$ is odd. Start with $v_2(D_0)=A$ and apply
\eqref{eq:diffrec} and Lemma~\ref{lem:two}; each step increases the valuation
by $t$.

For the 5-adic part, when $o=1$, $v_5(D_0)=v_5(a-1)=r$, and
\eqref{eq:orderLTE} applies immediately and at every subsequent step.
When $o=2$, $D_0$ is even and is not divisible by 5, so $y_1=r$; each later
step adds $r$. When $o=4$ and $a\equiv1\pmod4$, $4\mid D_0$ and again
$y_1=r$.

Finally suppose $o=4$ and $a\equiv3\pmod4$. Then $D_0\equiv2\pmod4$, so
$y_1=0$. Every odd-base tower $T_b$ with $b\geq1$ is congruent to $a$ modulo 4:
its exponent is odd. Consequently $4\mid D_b$ for $b\geq1$. Thus $y_2=r$,
and all subsequent steps add $r$. These four initial conditions give the four
lines of \eqref{eq:epsilonodd}.
\end{proof}

\begin{proposition}[Even bases]\label{prop:evenprofiles}
For even $a$ with $5\nmid a$, put $e=v_2(a)$. For all $b\geq1$,
\begin{equation}\label{eq:evenprofiles}
 x_b=eT_{b-1},\qquad y_b=\max\{0,(b+\varepsilon)r\},
\end{equation}
where
\begin{equation}\label{eq:epsiloneven}
\varepsilon=\begin{cases}
1,&o=1,\\
-1,&o=2,\\
-1,&o=4\text{ and }4\mid a,\\
-2,&o=4\text{ and }a\equiv2\pmod4.
\end{cases}
\end{equation}
Moreover, $x_b>y_b$ for every $b\geq2$.
\end{proposition}
\begin{proof}
The second factor in \eqref{eq:diffrec} is odd, so
$x_b=v_2(T_b)=eT_{b-1}$. If $o=1$, the 5-adic calculation starts with
$v_5(D_0)=r$. If $o=2$, $D_0$ is odd, so $y_1=0$; $D_1$ is even, so
$y_2=r$. For $o=4$, $D_0$ is odd, and
$D_1=T_2-a\equiv-a\pmod4$ because $4\mid T_2$. If $4\mid a$ this gives
$y_2=r$; otherwise $y_2=0$ and $y_3=r$, since $4\mid D_2$. Thereafter each
step adds $r$.

It remains to prove strict dominance, rather than merely eventual dominance.
For $o=1$, $a\geq5^r+1>3r$, so $x_2=ea>3r=y_2$. For $o=2$,
$a\geq5^r-1>r=y_2$. For $o=4$, the elementary inequality
$5^a>a^2+1$ for $a\geq2$ implies $r<a$, so $x_2>r\geq y_2$.
The inequality $5^a>a^2+1$ follows by induction from $a=2$.

Also $T_{j+1}=a^{T_j}\geq2T_j$, since $2^n\geq2n$ for every positive
integer $n$. Therefore $x_b\geq2^{b-2}x_2$. In the first case
$3\cdot2^{b-2}\geq b+1$; in the other cases $2^{b-2}\geq b-1$.
The displayed strict inequalities persist for all $b\geq2$.
\end{proof}

\begin{proposition}[Bases ending in 5]\label{prop:fiveprofiles}
For odd $a$ divisible by 5, put $f=v_5(a)$ and retain $A,t$ from
\eqref{eq:At}. Then
\begin{equation}\label{eq:fiveprofiles}
 x_b=A+bt,\qquad y_b=fT_{b-1}\qquad(b\geq1).
\end{equation}
If $a\geq15$, $y_b>x_b$ for all $b\geq2$. If $a=5$, this strict inequality
holds for $b\geq3$, and
\begin{equation}\label{eq:fiveexception}
 k_1=1,\quad k_2=5,\quad k_b=2b+2\ (b\geq3),\qquad B=4.
\end{equation}
\end{proposition}
\begin{proof}
The 2-adic calculation is unchanged. The factor $a^{D_{b-1}}-1$ in
\eqref{eq:diffrec} is prime to 5, giving the formula for $y_b$.
For $a\geq15$ we have $a>3t\geq A+2t$: for $t=2,3$ this is immediate;
for $t\geq4$ use $a\geq2^t-1>3t$. Thus $y_2=fa>x_2$.
The same doubling argument as above preserves strict dominance over the
linear function $x_b$. For $a=5$, direct substitution gives $k_1=1$, $k_2=5$, and
$y_3=3125>x_3=8$; doubling then preserves $y_b>x_b$ for all $b\geq3$,
giving the displayed $k_b$ values. The speeds are $1,4,3,2,2,\ldots$, proving $B=4$.
\end{proof}

\section{Why the minimal speed anchor is already safe}\label{sec:onset}

For nonunit bases, the preceding propositions immediately give exact onsets:
\begin{equation}\label{eq:Beven}
 B=\begin{cases}
2,&a\text{ even},\ o=1,\ e\geq2r,\\
3,&a\text{ even},\ o=1,\ e<2r,\\
3,&a\text{ even},\ o=4,\ a\equiv2\pmod4,\\
2,&a\text{ even in the remaining cases}.
\end{cases}
\end{equation}
Indeed, in the first two lines $k_1=\min(e,2r)$ and $k_b=(b+1)r$ for
$b\geq2$. The other lines follow from their zero-valued initial 5-adic
profiles. None has $B=1$.

For $5\mid a$, $a\geq15$, the onset is
\begin{equation}\label{eq:Bfive}
 B=\begin{cases}2,&k_1=A+t,\\3,&k_1<A+t,\end{cases}
 \qquad k_1=\min(A+t,f).
\end{equation}
Together with $B(5)=4$, these formulas place the anchor strictly inside the
2-adically dominant regime. No tie remains at the anchor.

For odd bases coprime to 5 we need to understand the minimum of two affine
functions. This small point is central to the argument.

\begin{lemma}[No isolated valuation tie at the anchor]\label{lem:anchor}
For the profiles in Proposition~\ref{prop:oddprofiles}, let $v=\min(t,r)$.
The speeds are eventually $v$. If $t\ne r$, the line with smaller slope is
\emph{strictly} lower at every integer $b\geq B$. If $t=r$, the two lines are
either strictly ordered everywhere or identical everywhere.
\end{lemma}
\begin{proof}
Suppose first $t>r$. Set
$\delta_b=x_b-y_b=A-\varepsilon r+b(t-r)$, which is strictly increasing.
For $b\geq2$, direct subtraction of the two possible minima shows
\begin{equation}\label{eq:criterion}
 V_b=r\quad\Longleftrightarrow\quad\delta_{b-1}\geq0.
\end{equation}
If $\delta_{b-1}<0$ and the smaller branch changes during the step, then
$V_b=r-\delta_{b-1}>r$; if it does not change, $V_b=t>r$.
Consequently $B\geq2$ implies $\delta_B>0$, and the inequality only strengthens.
If $B=1$, then $k_1=r$. A tie there would require
$x_1=A+t=r$, impossible since $A\geq1$ and $t>r$. So this case is strict too.

If $t<r$, exchange the two lines: set
$\delta_b=y_b-x_b=\varepsilon r-A+b(r-t)$. The analogue of
\eqref{eq:criterion} is $V_b=t$ if and only if $\delta_{b-1}\geq0$.
The same proof gives strictness when $B\geq2$. In fact $B=1$ cannot occur in
this case: if the lower-slope line is already minimal at height one, then
$k_1=A+t>t$; otherwise some later increment is not $t$.

Finally, if $t=r$ the difference $x_b-y_b=A-\varepsilon t$ is constant.
\end{proof}

Identical lines require $A=\varepsilon t$. Since $A>0$ and
$\varepsilon\in\{-1,0,1\}$, this is equivalent to
\begin{equation}\label{eq:balancecondition}
 \varepsilon=1,\qquad A=t=r.
\end{equation}
Thus $a\equiv1\pmod{20}$ and
$v_2(a-1)=v_5(a-1)=v\geq2$. Equivalently,
$a=1+10^v c$ with $\gcd(c,10)=1$. These are exactly the balanced bases.
For them $k_1=2v$ and $V_b=v$ for $b\geq2$, so $B=2$.

\subsection{Closed onset formula for odd units}\label{sec:closedB}
The proof gives a compact exact formula useful in computation. If $t=r$, then
$B=1$ when $k_1=t$, and $B=2$ otherwise. If $t\ne r$, define
\begin{equation}\label{eq:Cq}
 (C,q)=\begin{cases}(A-\varepsilon r,t-r),&t>r,\\
                         (\varepsilon r-A,r-t),&t<r.
          \end{cases}
\end{equation}
Here $q>0$ and $\delta_b=C+bq$. Then
\begin{equation}\label{eq:Bclosed}
 B=\begin{cases}
 1,&k_1=v\text{ and }C+q\geq0,\\
 \max\bigl\{2,1+\lceil-C/q\rceil\bigr\},&\text{otherwise}.
 \end{cases}
\end{equation}
This formula does not infer an onset from a finite run of equal speeds; it
proves that every later speed agrees.

\section{The phase dynamics in the three regimes}\label{sec:dynamics}

\subsection{The 2-adic regime}
If $x_b<y_b$, then $k_b=x_b$ and $D_b/10^{k_b}$ is odd and divisible by 5.
Its residue modulo 10 is therefore 5, and so is its negative. Hence $s_b=5$.
The anchor lemma and the nonunit calculations show that this conclusion holds
for \emph{all} $b\geq B$ whenever this is the eventual regime.

\subsection{The 5-adic regime: a four-element group}\label{sec:F5}
Suppose $y_b<x_b$ for all $b\geq B$. The base is prime to 5 and the constant
speed is $r$. Every $s_b$ is even and nonzero modulo 5. Let
\begin{equation}\label{eq:C5}
 u=\frac{a^o-1}{5^r},\qquad
 C_5=\begin{cases}a^a\pmod5,&a\text{ odd},\\1,&a\text{ even},\end{cases}
 \qquad \lambda=C_5u\,o^{-1}2^{-r}\in\F_5^\times.
\end{equation}
All inverses in this formula are taken in $\F_5$, so they exist.

For $b\geq B$, we have $o\mid D_b$. For odd bases this follows from
$4\mid D_b$ at every positive height; for even bases $B\geq2$ and
$4\mid D_b$ at those heights. Also $T_{b+1}=a^{T_b}\equiv C_5\pmod5$:
$T_b\equiv a\pmod4$ for odd bases and $T_b\equiv0\pmod4$ for even bases
at the relevant heights.

Apply Lemma~\ref{lem:five} to $a^{D_b}=(a^o)^{D_b/o}$ and use
$D_{b+1}=T_{b+1}(a^{D_b}-1)$. Since $k_{b+1}=k_b+r$, division by
$10^{k_b+r}$ gives
\begin{equation}\label{eq:geometric}
 \frac{D_{b+1}}{10^{k_{b+1}}}
 \equiv C_5u\,o^{-1}2^{-r}\frac{D_b}{10^{k_b}}\pmod5,
 \qquad s_{b+1}\equiv\lambda s_b\pmod5.
\end{equation}
The factor $2^{-r}$ is essential: the normalization removes powers of 10,
not just powers of 5.

Let $\ell:\F_5^\times\to\{2,4,6,8\}$ be the unique even representative
modulo 10; explicitly $\ell(1)=6$, $\ell(2)=2$, $\ell(3)=8$, and
$\ell(4)=4$. We have
\begin{equation}\label{eq:lift}
 s_{B+j}=\ell(s_B\lambda^j\bmod5).
\end{equation}
The multiplicative group $\F_5^\times$ has order 4. Multipliers $1,4,2,3$
have orders $1,2,4,4$, respectively. As the initial nonzero residue varies,
the resulting words are exactly the 16 members of $\Eset$ in
\eqref{eq:Eset}. This proves both the period bound and the complete
\emph{allowable} even-digit list. Actual witnesses are supplied below.

\subsection{Balanced valuations: the four exceptional words}

\begin{lemma}[A decimal first-order lifting identity]\label{lem:decimal}
Let $a=1+10^v c$, $v\geq2$, $\gcd(c,10)=1$. Then for every $b\geq0$,
\begin{equation}\label{eq:decimalunits}
 D_b\equiv10^{(b+1)v}c^{b+1}\pmod {10^{(b+1)v+1}}.
\end{equation}
\end{lemma}
\begin{proof}
The valuation identities already show
$v_2(D_b)=v_5(D_b)=(b+1)v$, including $b=0$.
The congruence starts with $D_0=10^v c$. Expand
\begin{equation}\label{eq:binomial}
 a^{D_b}-1=D_b10^v c+
     \sum_{j=2}^{D_b}\binom{D_b}{j}10^{jv}c^j.
\end{equation}
For $p=2,5$, the integer identity
$j\binom{D_b}{j}=D_b\binom{D_b-1}{j-1}$ gives
\[
 v_p\!\binom{D_b}{j}\geq v_p(D_b)-v_p(j).
\]
Every term with $j\geq2$ therefore has $p$-adic valuation at least
\[
 (b+2)v+(j-1)v-v_p(j)\geq(b+2)v+1.
\]
The last inequality holds because $v\geq2$ and $v_p(j)\leq j-1$.
Thus all nonlinear terms disappear modulo $10^{(b+2)v+1}$.
Multiply by $T_{b+1}\equiv1\pmod{10}$ in \eqref{eq:diffrec}, divide by
$10^{(b+2)v}$, and conclude that the normalized leading residue is multiplied
by $c$ at each step. Induction proves \eqref{eq:decimalunits}.
\end{proof}

Since $B=2$, the four possible words start with $-c^3$, not with $-c^2$.
They are
\begin{center}
\begin{tabular}{cccl}
\toprule
$c\bmod10$ & first four phase digits & reduced word & example base\\
\midrule
1 & $9,9,9,9$ & $9$ & $101$\\
3 & $3,9,7,1$ & $3971$ & $301$\\
7 & $7,9,3,1$ & $7931$ & $701$\\
9 & $1,9,1,9$ & $19$ & $901$\\
\bottomrule
\end{tabular}
\end{center}
The anchor, the sign convention, and the use of base 10 all matter in this table.
In particular, arbitrary rotations of these words need not be attainable at
$B$.

\begin{proof}[Proof of Theorem~\ref{thm:main}, except for the explicit witnesses]
Propositions~\ref{prop:oddprofiles}--\ref{prop:fiveprofiles} give eventual
positive speed. Section~\ref{sec:onset} proves that precisely one of strict
2-adic dominance, strict 5-adic dominance, or identical affine valuations
holds from the minimal anchor onward. The three calculations in this section
give the word $5$, one of the 16 words of $\Eset$, or one of the four words
of $\Uset$, respectively. In each case the phase period divides 4.
The witness table in Section~\ref{sec:witnesses} proves that none of the
21 possibilities is superfluous.
\end{proof}

\section{Why even bases allow only fourteen words}\label{sec:evenproof}
The finite-group argument alone permits all 16 even-digit words. The even-base
restriction comes from the initial residue and cannot be deduced just from
$B\leq3$.

For even $a$, $C_5=1$, so $\lambda=u\,o^{-1}2^{-r}$. The first phase residue
at the minimal anchor satisfies the following exact table, entirely in
$\F_5$:
\begin{center}
\begin{tabular}{@{}llcc@{}}
\toprule
Condition & additional condition & $B$ & $s_B\pmod5$\\
\midrule
$o=1$ & $e\geq2r$ & 2 & $-\lambda^3$\\
$o=1$ & $e<2r$ & 3 & $-\lambda^4$\\
$o=2$ & & 2 & $3\lambda$\\
$o=4$ & $4\mid a$, $a\equiv2\pmod5$ & 2 & $\lambda$\\
$o=4$ & $4\mid a$, $a\equiv3\pmod5$ & 2 & $2\lambda$\\
$o=4$ & $a\equiv2\pmod4$ & 3 & $3\lambda$\\
\bottomrule
\end{tabular}
\end{center}

To verify the first two rows, when $o=1$ the normalized 5-adic residue obeys
$D_b/5^{(b+1)r}\equiv u^{b+1}\pmod5$. Hence, once $k_b=(b+1)r$,
$s_b\equiv-\lambda^{b+1}$. Formula~\eqref{eq:Beven} gives the two values of
$B$.

In the next three rows, $B=2$, $T_2\equiv1\pmod5$, and
$D_1\equiv1-a\pmod5$. The lifting calculation gives
$s_2\equiv\lambda(a-1)$. This is $3\lambda$ when $a\equiv4\pmod5$,
$\lambda$ when $a\equiv2\pmod5$, and $2\lambda$ when $a\equiv3\pmod5$.
In the last row, $T_2\equiv-1\pmod5$ while $T_3\equiv1\pmod5$, so
$D_2\equiv2\pmod5$. Thus $s_3\equiv-2\lambda=3\lambda$.

For $\lambda=2$, the table permits initial residues $1,2,4$, but not 3.
The missing start gives exactly the word $8624$. For $\lambda=3$, it permits
$1,3,4$, but not 2; the missing start gives $2684$. For $\lambda=1$ all four
constant phases remain possible, and for $\lambda=4$ all four period-two
words remain possible. This proves the 14-word upper bound. The even-base
witnesses in the next section prove equality, completing
Theorem~\ref{thm:even}.

\section{Explicit witnesses and exact certificates}\label{sec:witnesses}
The following table gives a base for each of the 21 words and, when it exists,
an even base with the same word. The values $B$, $k_B$, and $s_B$ refer to the
base in the second column. The companion file \texttt{data/witnesses.csv}
also contains the two tower residues modulo $10^{k_B+1}$, which provide small
exact certificates for the first phase. Later phases follow from the proved
recurrence, not from extrapolation.

\begin{center}
\begin{tabular}{rrrrrrr}
\toprule
word & base $a$ & $B$ & $v$ & $k_B$ & $s_B$ & even witness\\
\midrule
2&9&1&1&1&2&48\\
4&11&2&1&3&4&36\\
5&5&4&2&10&5&---\\
6&52&2&1&1&6&52\\
8&2&3&1&1&8&2\\
9&101&2&2&6&9&---\\
19&901&2&2&6&1&---\\
28&18&3&2&2&2&18\\
46&3&2&1&1&4&12\\
64&16&2&1&3&6&16\\
82&39&1&1&1&8&88\\
2486&21&2&1&3&2&72\\
2684&7&2&2&2&2&---\\
3971&301&2&2&6&3&---\\
4268&6&3&1&4&4&6\\
4862&19&1&1&1&4&28\\
6248&4&2&1&1&6&4\\
6842&8&2&1&1&6&8\\
7931&701&2&2&6&7&---\\
8426&31&2&1&3&8&32\\
8624&77&1&1&1&8&---\\
\bottomrule
\end{tabular}
\end{center}

For example, base 7 has $o=4$, $r=2$, $A=1$, $t=3$, $B=2$, and $k_B=2$.
The residues $T_2\equiv543$ and $T_3\equiv343\pmod{1000}$ give $s_B=2$.
Furthermore $C_5=3$, $u\equiv1$, and $\lambda=3$, so
$2\mapsto6\mapsto8\mapsto4\mapsto2$. This independently yields $2684$.

For base 21, $T_2\equiv4421$ and $T_3\equiv2421\pmod{10000}$ give
$s_2=2$, and $\lambda=2$, hence $2486$. The significance of this example for
the currently listed OEIS values is discussed in Section~\ref{sec:audit}.

\section{The related 28-value conjecture}\label{sec:modular}
OEIS A376446 uses the prescribed anchor
\begin{equation}\label{eq:H}
 H=\nu(a)+2,\qquad
 \nu(a)=\begin{cases}r,&5\nmid a,\\t,&5\mid a,\end{cases}
\end{equation}
instead of $B$ \cite{oeis446}. Denote its reduced phase word by $\MPS(a)$.
The distinction is an actual rotation of a cycle, not a different notion of
phase.

\begin{lemma}[The prescribed anchor is safe]\label{lem:safeH}
For every admissible base, $B\leq H$.
\end{lemma}
\begin{proof}
Even bases have $B\leq3\leq r+2$. Bases ending in 5 have $B\leq t+2$ by
\eqref{eq:Bfive} and \eqref{eq:fiveexception}. For odd units with $t=r$,
$B\leq2$. If $t>r$, the increasing difference
$A-\varepsilon r+b(t-r)$ is nonnegative by $b=\max(1,r-1)$, so
\eqref{eq:Bclosed} gives $B\leq\max(2,r)\leq r+2$.

If $t<r$ and $\varepsilon=-1$, then $a\equiv3\pmod4$, so $A=1$ and the
difference is $-r-1+b(r-t)$; it is nonnegative by $b=r+1$.
If $\varepsilon=0$, it is $-A+b(r-t)$ with $A\leq t<r$, again nonnegative
by $r+1$. If $\varepsilon=1$, it is already positive at $b=0$, since
$r>A$. Formula~\eqref{eq:Bclosed} finishes the proof.
\end{proof}

\begin{theorem}[The 28-value modular classification]\label{thm:modular}
The exact range of $\MPS$ is
\begin{equation}\label{eq:MPSrange}
\Eset\cup\{5,9,19,91,1397,1793,3179,3971,7139,7931,9317,9713\}.
\end{equation}
This is the 28-value set conjectured in OEIS A376446.
\end{theorem}
\begin{proof}
By Lemma~\ref{lem:safeH}, strict 2-adic and 5-adic regimes still give $5$ and
words in $\Eset$. For a balanced base $a=1+10^v c$ we have $H=v+2$ and
\begin{equation}\label{eq:MPSbalanced}
 s_{H+j}\equiv-c^{v+3+j}\pmod {10}.
\end{equation}
If $c\equiv1$, this is always $9$. If $c\equiv9$, the parity of $v$ gives
$19$ or $91$. If $c\equiv3$ or 7, the four residues of $v$ modulo 4 give all
four rotations of the corresponding four-cycle. These account for exactly the
11 odd unit words in \eqref{eq:MPSrange}.

All of them are attained by $a=1+c10^v$ with $c\in\{1,3,7,9\}$ and
$2\leq v\leq5$. The words of $\Eset$, in the order shown in
\eqref{eq:Eset}, are attained by bases
\[
9,11,52,2,33,16,3,18,4,32,6,13,19,27,7,28,
\]
respectively; $a=5$ gives $5$. These are exact finite modular evaluations, also
recorded and checked in \texttt{data/modular\_witnesses.csv}.
\end{proof}

\section{A density consequence}
\begin{corollary}\label{cor:density}
Each of the four sets
\[
 \{a:\APS(a)=9\},\quad \{a:\APS(a)=19\},\quad
 \{a:\APS(a)=3971\},\quad\{a:\APS(a)=7931\}
\]
has natural density $1/900$ in the positive integers. Their union has density
$1/225$.
\end{corollary}
\begin{proof}
Theorem~\ref{thm:main} makes the balanced family necessary as well as
sufficient. For a specified residue $c_0\in\{1,3,7,9\}$ and a specified
$v\geq2$, the relevant bases form the single arithmetic progression
\[
 a\equiv1+c_0 10^v\pmod {10^{v+1}},
\]
with density $10^{-(v+1)}$. These progressions are disjoint as $v$ varies,
since $v=v_{10}(a-1)$ is determined by $a$. The union over $v\geq N$ is
contained in $a\equiv1\pmod{10^N}$, so its upper density is at most $10^{-N}$.
This bound justifies passage from finite sums to the infinite union. Therefore
\[
 \sum_{v=2}^\infty10^{-(v+1)}=\frac{1}{900}
\]
is the natural density for each specified word. Summing the four disjoint
families gives $1/225$.
\end{proof}

Thus the rarity of odd non-5 phase words has an exact arithmetic explanation.
The finite test through base $100000$ found 111 instances of each of these
four words. This count is a check on the formulas; the density proof does not
rely on it.

\section{Unbounded onset and a sharp family}\label{sec:sharp}
The phase alphabet is finite, but the minimal-speed onset is unbounded. A
uniform search through a fixed number of heights is therefore not a proof of
stabilization.

\begin{proposition}[A constructive sharp-onset family]\label{prop:sharp}
For every integer $r\geq3$, let $a_r$ be the least positive representative of the
Chinese remainder conditions
\begin{align}
 a_r&\equiv2^{r-1}-1\pmod {2^r},\label{eq:crt2}\\
 a_r&\equiv2^{5^{r-1}}\pmod {5^{r+1}}.\label{eq:crt5}
\end{align}
Then
\begin{equation}\label{eq:sharp}
 B(a_r)=r+2,\qquad v(a_r)=r-1,\qquad \APS(a_r)=5.
\end{equation}
Hence the general sufficient bound $B\leq r+2$ for odd units is sharp.
\end{proposition}
\begin{proof}
Condition \eqref{eq:crt2} gives $A=1$ and $t=r-1$. The other condition gives
$a_r\equiv2\pmod5$, so $o=4$, and $a_r\equiv3\pmod4$, so
$\varepsilon=-1$.

Put $n=5^{r-1}$, which is odd. Applying Lemma~\ref{lem:five} to
$(-4)^n-1=-(4^n+1)$ gives
\[
 v_5(4^{5^{r-1}}+1)=v_5((-4)-1)+(r-1)=r.
\]
The residue modulo $5^{r+1}$ in \eqref{eq:crt5} therefore ensures
$v_5(a_r^2+1)=r$ exactly. The two affine valuations are
\[
 x_b=1+b(r-1),\qquad y_b=(b-1)r,
 \qquad y_b-x_b=b-r-1.
\]
They are equal at $b=r+1$, and 2-adic dominance starts strictly at $b=r+2$.
The increment into the tied height is still $r$; the increment out is $r-1$.
Thus the first permanently constant speed is at $r+2$, where all subsequent
phases are 5.
\end{proof}

For example $r=3$ gives $a_3=1307$, $B=5$, and speed 2. One tested member
with $r=32$ is
\[
 a_{32}=311423787862411847003581666295807,
\]
whose onset is 34 and whose speed is 31. At that onset $k_B=1055$, so its first
phase can be independently checked using residues modulo $10^{1056}$, despite
the inaccessible size of the actual tower. The bound itself is known in the
congruence-speed literature \cite{ripaonnis}; the point here is the explicit
CRT construction and the exact explanation of its crossing height.

\section{Exact computation without constructing towers}\label{sec:algorithm}

\subsection{The main algorithm}
For a given admissible base, compute $A,t$ when $a$ is odd, and $o,r$ when
$5\nmid a$. Propositions~\ref{prop:oddprofiles}--\ref{prop:fiveprofiles} and
\eqref{eq:Beven}, \eqref{eq:Bfive}, \eqref{eq:Bclosed} give $B$, $v$, and $k_B$.
If the 2-adic valuation is smaller, return $5$. If they are equal at $B$, use
the balanced formula. Otherwise compute a single phase $s_B$ by modular
tetration modulo $10^{k_B+1}$ and propagate it with the multiplier
\eqref{eq:C5}. Four direct phase evaluations are included as an independent
internal check, rather than being needed by the mathematical algorithm.

A nonunit local valuation such as $v_2(a)T_{b-1}$ must not be constructed when
it is enormous. To compare it with a known integer $y$, it suffices to compute
$\min\{T_{b-1},\lfloor y/v_2(a)\rfloor+1\}$. This preserves the smaller
valuation and detects equality exactly.

\subsection{A reliable modular-tower routine}
The supplied \texttt{Tower.mod} handles moduli $m=2^u5^w$. For each prime power
$p^e\mid m$ there are two cases. If $p\nmid a$, reduce the exponent modulo
$\varphi(p^e)$ and recurse on the lower-height tower. If $p\mid a$, compute
the exponent only up to $\lceil e/v_p(a)\rceil$: at or above this threshold the
residue is zero; below it the exponent is small and known exactly. Combine
the two prime-power residues by the Chinese remainder theorem.

The recursion terminates because the height decreases. Every new modulus is
again 2,5-smooth. The exponent reduction in the unit case is justified by
Euler's theorem; in the nonunit case no unjustified exponent reduction is
made.

The base cases are $T_0=1$ and $\min\{T_h,1\}=1$.
For $h\geq1$ and $M>1$, to compute $\min\{T_h,M\}$, let $q$ be the least nonnegative integer with
$a^q\geq M$, compute $\min\{T_{h-1},q\}$ recursively, and return $M$ when
the threshold is reached. Otherwise evaluate the small exact power. All of
this uses integers, with no floating-point logarithms.

\begin{corollary}[Polynomial-time computability]\label{cor:complexity}
The word $\APS(a)$ can be computed in time polynomial in the bit length of $a$.
\end{corollary}
\begin{proof}
Write $L$ for that bit length. The valuations $A,t,r$ are $O(L)$, since they
are taken on $a\pm1$ or $a^o-1$ with $o\leq4$. The onset bound gives
$B=O(L)$ and the local profiles give $k_B=O(L^2)$. Thus the required modulus
has $O(L^2)$ bits, not a tower-sized representation.

In a memoized modular recursion a state is specified by a height at most
$B+1$ and the two exponents of its smooth modulus. Those exponents remain
$O(L^2)$; hence even the coarse bound on the number of possible states is
polynomial in $L$. Each state uses polynomial-size integer arithmetic,
modular exponentiation, and at most two recursive prime-power branches.
The capped-tower calculations also use polynomial-size integers and bounded
recursion depth. Standard multiplication and repeated squaring therefore
suffice for a polynomial-time algorithm. No optimal complexity bound is
claimed here.
\end{proof}

\subsection{A second algorithm used for checking}
The verification program also implements full-modulus Euler reduction with
exponent lifting. It computes $T_h\pmod m$ by reducing the exponent modulo
$\varphi(m)$ and adding $\varphi(m)$ when the actual exponent is at least that
large. If it is smaller, it uses the exact exponent instead.

For a prime power dividing $m$, an exponent at least $\varphi(m)$ is large
enough to force the nonunit residue to zero, since
$\varphi(m)\geq\varphi(p^e)\geq e$. In the unit factors, adding
$\varphi(m)$ does not change the residue. This algorithm does not split the
answer into prime-power residues and reconstruct it by CRT, so it exercises
a different reduction path. It shares only the capped-tower helper, which is
separately checked against directly constructed small towers.

\section{Reproducible verification and source audit}\label{sec:audit}

\subsection{What was actually checked}
The included program was executed with Python 3.13.5, using standard-library
integer arithmetic only. The saved run records:
\begin{center}
\begin{tabular}{lr}
\toprule
Check & count or range\\
\midrule
Admissible bases in $2\leq a\leq100000$ & $89999$\\
Post-onset phase positions checked & $719992$\\
Direct-integer residue, cap, and valuation checks & $3006$\\
Independent Euler-algorithm comparisons & $107$\\
Structured sharp-onset examples & $7$\\
Largest structured onset & $34$\\
Distinct minimal-anchor words & $21$\\
Distinct even-base minimal-anchor words & $14$\\
Distinct prescribed-anchor words & $28$\\
\bottomrule
\end{tabular}
\end{center}
For each sequentially tested base, the program checked eight post-onset
phases, the exact increment at each tested height, and the failure of the
limiting increment immediately before the proposed onset when $B>1$.
It also compared the four directly computed phases with the structural
prediction from one phase and its multiplier, or with the balanced formula.
The source code and its assertions make the scope of these tests explicit.

The sequential count excludes 1 and multiples of 10; it is therefore $89999$,
not $90000$. The elapsed time in the saved successful run was 28.344 seconds
on the execution environment used here. It is not a hardware-independent
performance guarantee. The JSON report and CSV files contain the exact
results, including the discrepancy described next.

\subsection{The value currently listed at base 21}
The accessed A376842 entry lists $6248$ at base 21 \cite{oeis842}. With the
height-one convention of A371074, the calculation is instead
\begin{equation}\label{eq:21calc}
 A=t=2,\quad r=1,\quad\varepsilon=1,\quad
 k_b=b+1\ (b\geq1),\quad V_1=2,\quad B=2.
\end{equation}
The onset $B(21)=2$ is independently the value listed in OEIS A372490
\cite{oeis2490}. From the exact residues
\[
 T_2\equiv4421,\qquad T_3\equiv2421\pmod {10000}
\]
and the multiplier $\lambda=2$ follows
\begin{equation}\label{eq:21answer}
 \APS(21)=2486.
\end{equation}
The phase at height one is 6, so $6248$ is the same cycle started one height
earlier. It is therefore consistent with an anchoring slip rather than a
failure of the periodicity or the 21-value range. The supplied checker records
this mismatch explicitly; it does not silently modify the source data.
Every other admissible listed value from base 2 through base 59 agrees with
the calculation.

\subsection{The height-zero ambiguity}
Definitions phrased uniformly as a difference between the numbers of stable
digits at heights $b$ and $b-1$ invite substituting $T_0=1$ at $b=1$.
The height-one convention recorded in A371074 instead counts the matching
digits of $a$ and $a^a$ directly. This is the convention used throughout this
paper, and it matches the examples $\APS(301)=3971$ and $\APS(901)=19$ in
Section 3 of Rip\`a's 2025 paper \cite{ripa2025}.

For a balanced base, substituting $v_{10}(a-1)=v$ as an initial count would
make $B=1$, not 2, and would give
\[
 9,\quad1397,\quad1793,\quad91
\]
for $c\equiv1,3,7,9$, respectively. These are different anchored invariants.
Thus $901$ is not a counterexample to the intended 21-value conjecture; it
only exposes why the convention must be fixed. We resolve this ambiguity
before stating the theorem rather than relying on an accidental rotation.

\subsection{The base-5 length convention}
Some discussions count only physically written digits at height two, assigning
four stable digits to $5^5=3125$ even though its congruence with the next tower
holds modulo $10^5$. Here we consistently use the exact modular count
$k_2=5$. Either choice has $B(5)=4$ and $\APS(5)=5$, so it does not affect the
classification. It does affect the short initial list of speeds and should
not be left implicit.

\subsection{Literature status and limits of the claim}
The target conjecture was checked in the original MathOverflow question, the
2025 arXiv PDF including its appendix, and the current OEIS entry. The PDF
appendix matters: the HTML rendering consulted does not display that final
appendix as completely as the PDF. The 2022 valuation paper was checked
separately to distinguish existing speed results from the phase problem.
A targeted search for the sequence identifier and phase-shift terminology
located no subsequent proof of the complete classification.

This evidence supports selecting the problem as an explicitly outstanding
conjecture in the consulted sources. It does not establish an exhaustive
bibliographic negative or guarantee novelty of every lemma and corollary.
No assertion of external peer review, Lean formalization, or Wolfram
verification is made. The calculations reported here were performed locally
in Python; the mathematical argument is supplied in full above.

\section{Interpretation and further questions}
The tower's size is largely irrelevant to its first changing decimal digit.
The relevant mechanism is a comparison of two local divisibility profiles.
When one profile is strictly smaller, normalization kills the contribution
from the other prime. What remains is either the single residue 5 or a
multiplicative orbit in the four-element group $\F_5^\times$. Exactly balanced
profiles retain a decimal unit, and the anchor fixes which rotation of its
orbit is seen.

The same reasoning explains why an arbitrary later height has the same first
changing digit relative to \emph{any} further tower. For $b\geq B$ and $h>b$,
\[
 T_h-T_b=\sum_{j=b}^{h-1}D_j\equiv D_b\pmod {10^{k_b+1}},
\]
because every $D_j$ with $j>b$ is divisible by $10^{k_b+v}$ and $v\geq1$.
Thus the phase really describes the first permanently unsettled digit, not
merely a coincidental difference between two isolated heights.

There are natural extensions that this article does not assert as proved.
One is a full density distribution over all 21 words, rather than the four
balanced ones. Another is a comparable classification in other radices,
where more local primes or nontrivial prime-power weights can change the
finite-state dynamics. A third is a machine-checked formalization of the
anchor lemma and the valuation-to-phase reduction. These are separate tasks;
they are not gaps in the decimal classification proved here.

\begin{resultbox}
\textbf{Conclusion.} Under the documented height-one convention, the
21-value conjecture and its 14-value even-base refinement follow from exact
local valuation formulas and a finite-group recurrence. The prescribed-height
28-value conjecture follows as a rotation classification. The most delicate
step is not computing enormous towers: it is proving that the first constant
speed occurs strictly after any isolated tie between the local valuations.
\end{resultbox}

\appendix
\section{Archive contents and reproduction}
The archive contains the complete \LaTeX{} source and compiled PDF, together
with the following files:
\begin{center}
\begin{tabular}{@{}p{0.44\linewidth}p{0.50\linewidth}@{}}
\toprule
File & purpose\\
\midrule
\texttt{code/phase\_tetration.py} & Exact smooth-modulus towers, valuations,
onsets, phase words, and the sharp CRT family.\\
\texttt{code/verify.py} & Sequential, direct-integer, independent-reduction,
structured, and published-data checks.\\
\texttt{data/verification.json} & Machine-readable results of the successful run.\\
\texttt{data/verification.txt} & Human-readable copy of that result.\\
\texttt{data/witnesses.csv} & All 21 witnesses, onsets, speeds, and exact
first-phase residue certificates.\\
\texttt{data/modular\_witnesses.csv} & Witnesses for all 28 prescribed-anchor words.\\
\texttt{data/sharp\_onsets.csv} & Structured examples including onset 34.\\
\texttt{README.md} & Commands, conventions, and the limits of the computational checks.\\
\texttt{SOURCE\_STATUS.md} & Sources, versions, and the initial-height/data audit.\\
\texttt{build.sh}, \texttt{build.ps1} & PDF build commands for common environments.\\
\bottomrule
\end{tabular}
\end{center}

From the extracted directory, use Python 3.9 or later:
\begin{verbatim}
python code/phase_tetration.py 57
python code/phase_tetration.py 901
python code/verify.py --max-base 100000
\end{verbatim}
Do not use Python's \texttt{-O} option, which disables assertions. No Python
packages beyond the standard library are required. The saved run used Python
3.13.5; the code targets Python 3.9+ syntax and APIs, but was not executed on
every intervening interpreter version.

Build the document with a \TeX{} installation containing \texttt{newpxtext},
\texttt{newpxmath}, \texttt{tcolorbox}, and the usual AMS packages:
\begin{verbatim}
pdflatex -interaction=nonstopmode -halt-on-error \
  tetration_phase_classification.tex
pdflatex -interaction=nonstopmode -halt-on-error \
  tetration_phase_classification.tex
\end{verbatim}
The bibliography is self-contained in the source, so Bib\TeX{} is not required.
No font files or copies of third-party papers are included.

\begin{thebibliography}{99}
\bibitem{ripa2020}
Marco Rip\`a, \emph{On the constant congruence speed of tetration},
Notes on Number Theory and Discrete Mathematics \textbf{26}(3) (2020), 245--260.
\href{https://doi.org/10.7546/nntdm.2020.26.3.245-260}{doi:10.7546/nntdm.2020.26.3.245-260}.

\bibitem{ripa2021}
Marco Rip\`a, \emph{The congruence speed formula},
Notes on Number Theory and Discrete Mathematics \textbf{27}(4) (2021), 43--61.
\href{https://doi.org/10.7546/nntdm.2021.27.4.43-61}{doi:10.7546/nntdm.2021.27.4.43-61}.

\bibitem{ripaonnis}
Marco Rip\`a and Luca Onnis, \emph{Number of stable digits of any integer tetration},
Notes on Number Theory and Discrete Mathematics \textbf{28}(3) (2022), 441--457.
\href{https://doi.org/10.7546/nntdm.2022.28.3.441-457}{doi:10.7546/nntdm.2022.28.3.441-457}.
Consulted preprint: \url{https://arxiv.org/pdf/2210.07956}.

\bibitem{ripa2025}
Marco Rip\`a, \emph{Graham's number stable digits: An exact solution},
Notes on Number Theory and Discrete Mathematics \textbf{31}(3) (2025), 607--616.
\href{https://doi.org/10.7546/nntdm.2025.31.3.607-616}{doi:10.7546/nntdm.2025.31.3.607-616}.
Consulted arXiv version 2, 12 September 2025:
\url{https://arxiv.org/pdf/2411.00015v2}; see Definition 3.2 and the appendix
on preprint page 9 for the target conjectures.

\bibitem{mo}
Marco Rip\`a, \emph{A new Conjecture at OEIS sequence A376842},
MathOverflow question 481090, 22 October 2024.
\url{https://mathoverflow.net/questions/481090/}.
Accessed 19 September 2026; the retrieved question displayed no answer.

\bibitem{oeis842}
OEIS Foundation Inc., \emph{A376842}, entry by Marco Rip\`a.
\url{https://oeis.org/A376842/internal}.
Accessed 19 September 2026; revision 39, 30 March 2026.

\bibitem{oeis446}
OEIS Foundation Inc., \emph{A376446}, entry by Marco Rip\`a.
\url{https://oeis.org/A376446/internal}.
Accessed 19 September 2026; revision 40, 13 December 2025.

\bibitem{oeis1074}
OEIS Foundation Inc., \emph{A371074}, entry by Marco Rip\`a.
\url{https://oeis.org/A371074}.
Accessed 19 September 2026. The comments explicitly identify the total
height-one matching-digit count with the height-one speed.

\bibitem{oeis2490}
OEIS Foundation Inc., \emph{A372490}, entry by Marco Rip\`a.
\url{https://oeis.org/A372490}; table
\url{https://oeis.org/A372490/b372490.txt}.
Accessed 19 September 2026; used to cross-check the onset at base 21.
\end{thebibliography}
\end{document}
